\section{Universal parity identities beyond commensurate scales}
\label{clc:sec:parity}

At negative integer total orders the resolvent power becomes a
polynomial. A local residue at zero then yields exact identities
without a common period. Put
\begin{equation}\label{clc:eq:Bq}
 B_{\bq}(z)=z^rG_{\bq}(z)
       =\prod_{j=1}^r\frac{\pi z/q_j}{\sin(\pi z/q_j)}.
\end{equation}
This is an even analytic function near zero with $B_{\bq}(0)=1$.

\begin{theorem}[Reflection-residue parity law]
\label{clc:thm:parity}
Let $q_j>0$ be arbitrary and let $N\ge0$ be an integer. Then
\begin{align}
 &\CC_{\bq,A}(-N;0)-(-1)^{r+N}\CC_{\bq,-A}(-N;0)\notag\\
 &\hspace{18mm}=[z^{r-1}](z+A)^N B_{\bq}(z).
 \label{clc:eq:parity}
\end{align}
The two ordinary integrals are covered simultaneously when
$|\Rea A|<q_*$. At these integer orders each side is a polynomial in
$A$, so the identity also gives their polynomial continuation to all $A$.
In particular, if $r+N$ is odd,
\begin{equation}\label{clc:eq:parity-central}
 \boxed{\displaystyle
 \CC_{\bq,0}(-N;0)=\frac12[z^{r-N-1}]B_{\bq}(z).}
\end{equation}
A coefficient with negative index is zero.
\end{theorem}
\begin{proof}
At $W=-N$, the contour integrand is
$G_{\bq}(p)(p+A)^N$. Choose $0<c<q_*$. At integer order the expression
is a polynomial in $A$, so the contour can be considered independently
of any branch of a resolvent power. Move the vertical contour from $c$
to $-c$. The only pole crossed is zero; horizontal edges vanish by
exponential decay. If $I_c(A)$ denotes the integral on $\Rea p=c$,
\[
 I_c(A)-I_{-c}(A)=\Res_{p=0}G_{\bq}(p)(p+A)^N.
\]
Since $G_{\bq}(-p)=(-1)^rG_{\bq}(p)$, changing variable $p\mapsto-p$
in the second vertical integral gives
$I_{-c}(A)=(-1)^{r+N}I_c(-A)$.
The residue on the right is exactly
$[z^{r-1}](z+A)^NB_{\bq}(z)$. On the common ordinary-integral domain,
$I_c(\pm A)=\CC_{\bq,\pm A}(-N;0)$, proving
\eqref{clc:eq:parity}; the polynomial continuation follows at once.
At $A=0$ and $r+N$ odd, the left side is twice the desired value,
giving \eqref{clc:eq:parity-central}.
\end{proof}

\subsection{Odd-depth logistic integrals}
For $r=2h+1$ and $N=0$, the theorem gives
\begin{equation}\label{clc:eq:odd-depth}
 K_{\bq}(0)=\frac12[z^{2h}]
                   \prod_{j=1}^{2h+1}\frac{\pi z/q_j}{\sin(\pi z/q_j)}.
\end{equation}
Every such value is a homogeneous polynomial in the inverse squared
scales with rational coefficients times $\pi^{2h}$.
To see both the rationality and the exact construction, use
\[
 \frac{u}{\sin u}=1+\frac{u^2}{6}+\frac{7u^4}{360}+\cdots,
\]
whose coefficients are rational by formal division of the sine series.
The first two nontrivial cases are
\begin{align}
 K_{(q_1,q_2,q_3)}(0)
   &=\frac{\pi^2}{12}\sum_{j=1}^3q_j^{-2},
 \label{clc:eq:universal3}\\
 K_{(q_1,\ldots,q_5)}(0)
   &=\frac{\pi^4}{720}
       \left(7\sum_{j=1}^5q_j^{-4}
               +10\sum_{i<j}q_i^{-2}q_j^{-2}\right).
 \label{clc:eq:universal5}
\end{align}
Equation \eqref{clc:eq:universal3} is the ordinary double integral
\eqref{clc:eq:intro-cubic}. For $(q_1,q_2,q_3)=(1,\sqrt2,\sqrt3)$,
for instance, its value is $11\pi^2/72$, despite the absence of a
common scale period. For $(1,2,3)$ it recovers $49\pi^2/432$ from the
explicit Laurent table.

More generally,
\begin{equation}\label{clc:eq:negative-examples}
 \CC_{(q_1,q_2),0}(-1;0)=\frac12,\qquad
 \CC_{(q_1,q_2,q_3),0}(-2;0)=\frac12.
\end{equation}
If $r+N$ is odd and $N>r-1$, the value is exactly zero.
Such identities involve negative-order Lerch or polylogarithmic
profiles, which can change sign, and so do not contradict positivity
of the zero-order logistic convolution.

At nonzero $A$, equation \eqref{clc:eq:parity} supplies a finite
reflection identity instead of an evaluation of each side separately.
For example, at $r=2$, $N=2$ the residue is $2A$, so
\[
 \CC_{(q_1,q_2),A}(-2;0)-\CC_{(q_1,q_2),-A}(-2;0)=2A.
\]
This exact reflection is independent of the two scales. The theorem
makes no assertion that the complementary-parity central values have
a comparable universal polynomial form.
